\label{chap:83-13}

\section{Preconstitutive construction of the bidirectional action response and the vacuum operator}
\label{p12-83-c13-s1:ch:vacuum}

\subsection{Derivation of the action scale from the dodecaphasic state}

The action decade is determined intrinsically, without adopting it as an
inherited unit. The local determinantal functional is applied to the previously
constructed HMT output \(\alpha_{\HMT}\) and defines
\begin{equation}
\Sigma_H=\varphi_{\HMT}\alpha_{\HMT}^{16}.
\label{p12-83-c13-s1:eq:action-determinantal-reader}
\end{equation}
The exponent \(16\) is the order of the local cokernel of the signed closure; its decimal valuation fixes
\begin{equation}
\operatorname{ord}_{10}(\Sigma_H)=-34.
\label{p12-83-c13-s1:eq:action-decade}
\end{equation}

The origin of the exponent can be verified within this volume. On the
local orbit of four leaves, the Hadamard block acts
\begin{equation}
H_4=
\begin{pmatrix}
1&1&1&1\\
1&1&-1&-1\\
1&-1&1&-1\\
1&-1&-1&1
\end{pmatrix}.
\label{p12-83-c13-s1:eq:action-hadamard}
\end{equation}

\begin{lemma}[Local index of the action line]
The form normal of Smith of \(H_4\) is
\begin{equation}
\operatorname{SNF}(H_4)=\operatorname{diag}(1,2,2,4).
\label{p12-83-c13-s1:eq:action-smith}
\end{equation}
Consequently,
\[
\operatorname{coker}(H_4)\simeq
\Z/2\Z\oplus\Z/2\Z\oplus\Z/4\Z,
\qquad
|\operatorname{coker}(H_4)|=16.
\]
\end{lemma}

\begin{proof}
The greatest common divisor of the entries is one; that of the nonzero
\(2\times2\) minors is two; that of the \(3\times3\) minors is four; and
\(|\det H_4|=16\). The successive quotients of these determinantal divisors
are \(1,2,2,4\), which prove \cref{p12-83-c13-s1:eq:action-smith}. The product of the
invariants gives the order of the cokernel.
\end{proof}

The norm of the section \(\alpha_{\HMT}\) over those sixteen sheets
produces \(\alpha_{\HMT}^{16}\); the self-scaling \(\varphi_{\HMT}\) acts
only once on the basis, not once per sheet. Hence
\cref{p12-83-c13-s1:eq:action-determinantal-reader}. The certified comparisons
\begin{equation}
\alpha_{\HMT}^{16}<10^{-34}
<\varphi_{\HMT}\alpha_{\HMT}^{16}<10^{-33}
\label{p12-83-c13-s1:eq:action-decade-bounds}
\end{equation}
prove, without introducing the SI unit, that the decimal order is
exactly \(-34\). The terminal evaluation begins with
\begin{equation}
\Sigma_H
=1.046243838104756580184229208303089\ldots\times10^{-34}.
\label{p12-83-c13-s1:eq:action-sigma-value}
\end{equation}
The substitution of the kernel by three copies with exponent \(16^3\), or the repeated application of \(\varphi\) on each sheet, would alter the genealogy. Both operations constitute falsification criteria rather than equivalent conventions. The causal chain that is preserved is therefore,
\[
\APP\to\TRIT\to\TPK\to U_{12}^{\rm sign}
\to\alpha_{\HMT}\to\Sigma_H
\to10^{-34}\U_S\to
\{\hbar_{\rm pre},\hbar_{\rm ret}\}.
\]
This volume does not rederive \(\alpha_{\HMT}\), but develops the metrological geometry
subsequent to the selection of the type and decade by
\(\Sigma_H\). The integer \(-34\) is neither a mantissa nor an adjustment to the International
System, but the terminal valuation of the local functional.

\subsection{Oriented sections of the action line}

The realizations before and after the return constitute two sections
of the same action line:

\begin{equation}
\hbar_{\rm pre}=H_5\,10^{-34}\U_S,\qquad
\hbar_{\rm ret}=\eta_{\rm ret}\,10^{-34}\U_S.
\label{p12-83-c13-s1:eq:action-twoway}
\end{equation}

The common decade determines the line, while the oriented
information is expressed through

\begin{equation}
R_{\rm act}=\frac{\hbar_{\rm pre}}{\hbar_{\rm ret}}
=\frac{H_5}{\eta_{\rm ret}},
\qquad
\xi_{\rm act}=\frac12\log R_{\rm act}.
\label{p12-83-c13-s1:eq:action-rapidity}
\end{equation}

The first quantity provides a multiplicative coordinate and the second
an additive coordinate of rapidity type. They do not represent two independent
Planck constants, but the two orientations of the same section
with respect to the return.

If \(h_a=2\pi\hbar_a\), \(a\in\{\mathrm{pre},\mathrm{ret}\}\), a quantity proportional to \(h_a^{1/2}\) changes as \(R_{\rm act}^{1/2}\), one proportional to \(h_a^{-1/2}\) changes as \(R_{\rm act}^{-1/2}\), and a quotient in which the action cancels is a bidirectional invariant.

\subsection{Bilayer spectral decomposition of the angular operator}

Let \(R=R^*=R^{-1}\) be the active involution and

\[
P_\pm=\frac{I\pm R}{2},
\qquad
T=q_+P_+ +q_-P_-,
\qquad 0<q_+<q_-<1.
\]

With
\[
x=A_{\rm rad}=\frac{\pi A}{180},\qquad
y=C^*_{\rm rad}=\frac{\pi C^*}{180},
\]
the two channels are constructed before evaluating the vacuum:
\begin{equation}
q_+=e^{-x-y},\qquad q_-=e^{-x+y}.
\label{p12-83-c13-s1:eq:vacuum-angular-channels}
\end{equation}
The exchange of sheets reverses \(C^*\), that is, \(y\mapsto-y\), and permutes \(q_+\leftrightarrow q_-\). The pair is thus the spectral reading of the same angular antecedent, not two parameters fitted separately.

\begin{definition}[Constitutive response]
\begin{equation}
\mathcal V_{90,120}=(I-T^{90})(I-T^{120})^{-1}.
\label{p12-83-c13-s1:eq:vacuum-operator}
\end{equation}
\end{definition}

Since \(\|T\|<1\), the denominator is positive and invertible. Functional calculus gives

\begin{equation}
\mathcal V_{90,120}=r_+P_+ +r_-P_-,
\qquad
r_\pm=\frac{1-q_\pm^{90}}{1-q_\pm^{120}}.
\label{p12-83-c13-s1:eq:vacuum-eigenvalues}
\end{equation}

The assertion of invertibility deserves to be made explicit. For each leaf,
\(0<q_\pm^{120}<1\); therefore
\[
I-T^{120}=(1-q_+^{120})P_++(1-q_-^{120})P_-
\]
is strictly positive and
\[
(I-T^{120})^{-1}
=\frac{P_+}{1-q_+^{120}}+
 \frac{P_-}{1-q_-^{120}}.
\]
The response \(\mathcal V_{90,120}\) is not obtained by solving four
separate scalar equations: it is the result of the functional calculus of a
single positive contraction. This observation prevents permeability,
permittivity, impedance and propagation from losing their shared antecedent.

\begin{theorem}[Joint spectral determination of the four constitutive relations]
The four constituent characters
\begin{equation}
\widehat\varepsilon=r_+^2,\qquad
\widehat\mu=r_-^2,\qquad
\widehat Z=\frac{r_-}{r_+},\qquad
\widehat c=\frac1{r_+r_-}
\label{p12-83-c13-s1:eq:vacuum-four}
\end{equation}
are spectral functions of \(\mathcal V_{90,120}\). They satisfy
\begin{equation}
\widehat\mu\widehat\varepsilon=\widehat c^{-2},
\qquad
\frac{\widehat\mu}{\widehat\varepsilon}=\widehat Z^2.
\label{p12-83-c13-s1:eq:vacuum-relations}
\end{equation}
\end{theorem}

\begin{proof}
The identities are obtained by substituting \cref{p12-83-c13-s1:eq:vacuum-four}. No dimensional chart nor any exterior value intervenes.
\end{proof}

\begin{corollary}[Reconstruction of the two sheets]
The four characters do not contain four degrees of freedom. Any of
the pairs \((\widehat\mu,\widehat\varepsilon)\),
\((\widehat Z,\widehat c)\), or \((\mathcal E,\mathcal B)\) reconstructs the
two positive eigenvalues:
\begin{align}
r_-&=\sqrt{\widehat\mu}
=\sqrt{\widehat Z/\widehat c},
&r_+&=\sqrt{\widehat\varepsilon}
=\frac1{\sqrt{\widehat Z\widehat c}},
\label{p12-83-c13-s1:eq:vacuum-reconstruct-r}\\
\log r_-&=\frac{\mathcal E+\mathcal B}{2},
&\log r_+&=\frac{\mathcal E-\mathcal B}{2}.
\label{p12-83-c13-s1:eq:vacuum-reconstruct-log}
\end{align}
In particular, orientation is lost if only the product is retained,
but is recovered by adding the quotient.
\end{corollary}

\begin{proof}
The first line is the positive root of the identities
\cref{p12-83-c13-s1:eq:vacuum-four,p12-83-c13-s1:eq:vacuum-relations}. The second inverts the linear
system
\(\mathcal E=\log r_-+\log r_+\),
\(\mathcal B=\log r_--\log r_+\).
\end{proof}

This corollary expresses a central property of HMT: product and quotient are
complementary functionals of the same state. The product determines the
propagation in common, whereas the quotient preserves the relative orientation
of the inner and outer sheets.

In the V2 certificate,

\begin{align*}
r_+&=0.9999996285482493631786952281\ldots,\\
r_-&=0.9997241371895183641315165853\ldots,\\
\widehat\varepsilon&=0.9999992570966367027604416155\ldots,\\
\widehat\mu&=0.9994483504793269350899815559\ldots,\\
\widehat Z&=0.9997245085389372154555543466\ldots,\\
\widehat c&=1.0002763104861575261063862689\ldots.
\end{align*}

These numerals are the terminal recognition of \cref{p12-83-c13-s1:eq:vacuum-operator}; they do not define its exponents \(90,120\).

\subsection{Harmonic factorization of sectors \texorpdfstring{\(90\)}{90} and \texorpdfstring{\(120\)}{120}}

The harmonic semigroup is

\[
\langle90,120\rangle=30\langle3,4\rangle.
\]

If \(s=q^{30}\), then

\begin{equation}
\frac{1-q^{90}}{1-q^{120}}
=\frac{1-s^3}{1-s^4}
=\frac{1+s+s^2}{1+s+s^2+s^3}.
\label{p12-83-c13-s1:eq:three-four-completion}
\end{equation}

The identity also separates the response from its defect. If
\[
F_{3\to4}(s)=\frac{1+s+s^2}{1+s+s^2+s^3},
\qquad 0<s<1,
\]
then
\begin{equation}
d_{3\to4}(s):=1-F_{3\to4}(s)
=\frac{s^3}{1+s+s^2+s^3}>0.
\label{p12-83-c13-s1:eq:vacuum-defect34}
\end{equation}
The additional term is preserved exactly as a positive
defect. Moreover,
\begin{equation}
F_{3\to4}'(s)
=-\frac{s^2(3+2s+s^2)}{(1+s+s^2+s^3)^2}<0.
\label{p12-83-c13-s1:eq:vacuum-monotonic34}
\end{equation}
Consequently, the order of the angular channels uniquely determines
the order of the constitutive responses. The assignment of
\(\widehat\mu\) and \(\widehat\varepsilon\) to the respective sheets is obtained
analytically, without resorting to their decimal expansions.

Each sheet of the vacuum constitutes the exact completion from three to four
terms of the same elementary mode \(30\). Consequently, permeability
and permittivity are interpreted as two quadratic functionals of the same
completion evaluated on conjugate sheets, and not as independent
constitutive parameters.

Since the function \((1-q^{90})/(1-q^{120})\) decreases in \(0<q<1\) and \(q_+<q_-\), we obtain

\[
0<r_-<r_+<1,\qquad
0<\widehat\mu<\widehat\varepsilon<1,\qquad
0<\widehat Z<1<\widehat c.
\]

The limits of the functional determine its geometry. When \(q\to0^+\),
\(F(q)\to1\): the four readings approach unity and the memory
defect disappears. When \(q\to1^-\),
\[
F(q)\longrightarrow\frac{90}{120}=\frac34.
\]
The response therefore takes values in the positive interval \((3/4,1)\). The
endpoint \(3/4\) is not a fitted constant: it is the quotient of the two
channel lengths and the continuous realization of the
three--to--four completion.

\subsection{Sheet involution and preservation of orientation}

The sheet exchange \(C^*\mapsto-C^*\) acts by

\begin{equation}
\widehat\mu\longleftrightarrow\widehat\varepsilon,\qquad
\widehat Z\longmapsto\widehat Z^{-1},
\qquad
\widehat c\longmapsto\widehat c.
\label{p12-83-c13-s1:eq:vacuum-sheet}
\end{equation}

The common propagation is even; impedance preserves orientation.
This distinction will reappear in Barbero and CKM: not all functionals
associated with the change of sheet have the same parity.

\subsection{Logarithmic characters and tracial compatibility}

Define

\[
\mathcal E=\log(r_+r_-),\qquad
\mathcal B=\log(r_-/r_+).
\]

With the normalized trace \(\tau=\Tr/2\),

\begin{equation}
2\tau(\log\mathcal V)=\mathcal E,
\qquad
-2\tau(R\log\mathcal V)=\mathcal B.
\label{p12-83-c13-s1:eq:vacuum-traces}
\end{equation}

The pair \((\mathcal E,\mathcal B)\) constitutes a linear coordinate system for the operator
logarithm:
\begin{equation}
\log\mathcal V
=\frac{\mathcal E}{2}I-\frac{\mathcal B}{2}R.
\label{p12-83-c13-s1:eq:vacuum-log-decomposition}
\end{equation}
The scalar component is invariant under sheet exchange; the
oriented component changes sign. This decomposition directly
demonstrates the parity of propagation and the oriented character of
impedance.

The nonadic amplification

\[
\mathcal V_m=\mathcal V\otimes I_{9^m}
\]

is compatible with the expectations \(E_m=\mathrm{id}\otimes\tau_9\):

\[
E_m(\mathcal V_{m+1})=\mathcal V_m.
\]

Therefore, \(\mathcal E\) and \(\mathcal B\) are not peculiarities of a \(2\times2\) matrix; they persist in the UHF tower and in its tracial closure \(II_1\).

More precisely, for every polynomial \(p\) and afterwards, by continuity,
for every continuous function on the spectrum,
\begin{equation}
E_m\bigl(p(\mathcal V_{m+1})\bigr)=p(\mathcal V_m).
\label{p12-83-c13-s1:eq:vacuum-functional-refinement}
\end{equation}
Compatibility does not preserve only two numbers: it preserves the entire
functional algebra generated by the response. In particular, it preserves its
spectral projectors, powers, logarithm, and quadratic forms
that measure the defects \cref{p12-83-c13-s1:eq:vacuum-defect34}.

\subsection{Dimensional realization of the constitutive relations}

Dimensional Mechanics determines

\begin{align}
Z_{\rm vac}&=\widehat Z\,u_Su_Q^{-2},\\
c_{\rm vac}&=\widehat c\,u_C,\\
\mu_{\rm vac}&=\widehat\mu\,u_Su_C^{-1}u_Q^{-2},\\
\varepsilon_{\rm vac}&=\widehat\varepsilon\,u_S^{-1}u_C^{-1}u_Q^2.
\label{p12-83-c13-s1:eq:vacuum-sections}
\end{align}

The constitutive identities hold on the straight lines:

\begin{equation}
\mu_{\rm vac}\varepsilon_{\rm vac}=c_{\rm vac}^{-2},
\qquad
\mu_{\rm vac}/\varepsilon_{\rm vac}=Z_{\rm vac}^2.
\end{equation}

The SI chart comes later. It is not necessary to declare \(c\) primitive and then solve \(\mu\varepsilon=c^{-2}\): the three objects descend jointly from \(\mathcal V\) and from the declared lines.

The separation of levels is expressed by the diagram
\[
\begin{array}{c}
(A,C^*)\longrightarrow T\longrightarrow\mathcal V
\longrightarrow
(\widehat\mu,\widehat\varepsilon,\widehat Z,\widehat c)
\\[2mm]
\Big\downarrow\text{ line charts}\hspace{36mm}
\Big\downarrow\text{ final evaluation}
\\[2mm]
(u_S,u_Q,u_C)\longrightarrow
(\mu_{\rm vac},\varepsilon_{\rm vac},Z_{\rm vac},c_{\rm vac})
\longrightarrow\text{ SI coordinates}.
\end{array}
\]
The first line is dimensionless and operatorial; the second determines types of
quantity; the third selects a metrological chart. Altering the final chart
changes neither the eigenvalues nor the constitutive identities. In this
precise sense, mathematical--physical unification consists in distinguishing
the genealogical invariant from its coordinate system, not in removing the
dimensional structure.

\begin{proposition}[Spectral minimality]
Let \(W\) be a positive two-sheet operator with simple spectrum
\(\{r_+,r_-\}\). If four positive characters satisfy
\(\mu\varepsilon=c^{-2}\) and \(\mu/\varepsilon=Z^2\), if their propagation in common is determined by \(c^{-1}=\det W=r_-r_+\), and if its orientation
fixes \(Z=r_-/r_+\), then necessarily
\[
(\mu,\varepsilon,Z,c)
=\left(r_-^2,r_+^2,\frac{r_-}{r_+},\frac1{r_-r_+}\right).
\]
Therefore, \cref{p12-83-c13-s1:eq:vacuum-four} is not one parametrization among many:
it is the minimal positive realization compatible with product, quotient and
orientation.
\end{proposition}

\begin{proof}
From \(Z=r_-/r_+\) and \(c^{-1}=\det W=r_-r_+\) the positive roots are obtained uniquely
roots. The two constitutive identities then force afterwards
\(\mu=r_-^2\) and \(\varepsilon=r_+^2\).
\end{proof}

\subsection{Angular determinant and Weinberg electroweak relation}

The determinant of \(T\) is

\begin{equation}
\det T=q_+q_-=D_A.
\label{p12-83-c13-s1:eq:detT}
\end{equation}

The constitutive sector of the vacuum evaluates \(F(T)^2\), where
\(F(z)=(1-z^{90})/(1-z^{120})\), whereas the electroweak angular
sector evaluates \(\det T\). These are two functions of a
common antecedent, not a horizontal factorization between constants.
This distinction is developed in the \cref{ch:ew}.

\begin{falsifier}
The block is refuted if any of these exact identities is violated:
positivity of \(I-T^{120}\), equality
\cref{p12-83-c13-s1:eq:three-four-completion}, relations \cref{p12-83-c13-s1:eq:vacuum-relations}, or
compatibility \(E_m(\mathcal V_{m+1})=\mathcal V_m\). An SI coincidence
cannot compensate for an operator failure.
\end{falsifier}

\begin{statusbox}
Operator, spectrum, completion \(3\to4\), parity, amplification, and constitutive sections form an exact internal theorem. Their assembly into a single operator is established in this formulation; the angular data and the lines, previously derived in HMT, constitute the explicit starting structure.
\end{statusbox}


\section{Quantum--electric constants, thermal scale, and Planck system}
\label{p12-83-c13-s2:ch:metrology}

\subsection{Derivation of the charge section from the vacuum}

With \(Z_{\rm vac}\), \(h_a=2\pi\hbar_a\) and \(\alpha_{\HMT}\) fixed, the charge of sheet \(a\in\{\mathrm{pre},\mathrm{ret}\}\) is

\begin{equation}
e_a=\sqrt{\frac{2\alpha_{\HMT}h_a}{Z_{\rm vac}}}.
\label{p12-83-c13-s2:eq:charge}
\end{equation}

The section \(e_a\) is not introduced through an International System
coordinate intended to reconstruct \(\alpha\). The causal order proceeds from
\(\alpha_{\HMT}\), the action, and the vacuum toward the charge section.

\subsection{Von Klitzing and Josephson constants, conductance and magnetic flux}

From \cref{p12-83-c13-s2:eq:charge} we deduce

\begin{align}
R_K&=\frac{h_a}{e_a^2}=\frac{Z_{\rm vac}}{2\alpha_{\HMT}},
\label{p12-83-c13-s2:eq:rk}\\
G_0&=\frac{2e_a^2}{h_a}=\frac{4\alpha_{\HMT}}{Z_{\rm vac}},
\label{p12-83-c13-s2:eq:g0}\\
\Phi_{0,a}&=\frac{h_a}{2e_a}
=\sqrt{\frac{h_aZ_{\rm vac}}{8\alpha_{\HMT}}},
\label{p12-83-c13-s2:eq:phi0}\\
K_{J,a}&=\frac{2e_a}{h_a}=\Phi_{0,a}^{-1}.
\label{p12-83-c13-s2:eq:kj}
\end{align}

Consequently,

\begin{equation}
R_KG_0=2,\qquad G_0Z_{\rm vac}=4\alpha_{\HMT},
\qquad K_{J,a}\Phi_{0,a}=1.
\label{p12-83-c13-s2:eq:quantum-electrical-identities}
\end{equation}

The bidirectional covariance takes the form

\[
R_K^{\rm pre}=R_K^{\rm ret},\qquad
G_0^{\rm pre}=G_0^{\rm ret},
\]

\[
\frac{\Phi_{0,\rm pre}}{\Phi_{0,\rm ret}}=R_{\rm act}^{1/2},
\qquad
\frac{K_{J,\rm pre}}{K_{J,\rm ret}}=R_{\rm act}^{-1/2}.
\]

The resistance and conductance relations are independent of the
orientation of action. The flux--frequency pair preserves that
orientation and transforms its two components reciprocally.

\subsection{Relation thermal intrinsic of the scale \texorpdfstring{\(108\)}{108}}

The chronological structure \(108\) fixes

\begin{equation}
k_B\Theta_{\rm clk}\log3
=\frac{h_{\rm int}}{108t_0}
=\frac{2\pi\hbar_{\rm int}}{108t_0}.
\label{p12-83-c13-s2:eq:boltzmann-clock}
\end{equation}

This equality does not separately determine a numerical coordinate of \(k_B\) and another of \(\Theta_{\rm clk}\) without an additional normalization. It does determine their product as a section of energy. The factor \(\log3\) has three compatible generators:

\begin{equation}
\log3=C_1+C_2-2C_0
=\frac{S_{\TRIT}}{k_B}
=\frac{E_P}{k_B\Theta_{\rm clk}}.
\label{p12-83-c13-s2:eq:log3-triple}
\end{equation}

The three equalities correspond to residual, informational and
thermodynamic realizations of the same scalar. The coincidence of value does not identify
their respective domains.

\subsection{Planck System and Bidirectional Covariance}

The CT108 selector constructs

\begin{equation}
\theta_{108}=\left(\frac{54}{\pi}\right)^2,
\qquad
\ell_P=\frac{54}{\pi}\ell_0,
\qquad
t_P=\frac{54}{\pi}t_0,
\label{p12-83-c13-s2:eq:planck-length-time}
\end{equation}

and

\begin{equation}
E_P=\frac{\hbar_{\rm int}}{t_P}
=\frac{\pi\hbar_{\rm int}}{54t_0}.
\label{p12-83-c13-s2:eq:planck-energy}
\end{equation}

Comparing \cref{p12-83-c13-s2:eq:boltzmann-clock,p12-83-c13-s2:eq:planck-energy} gives

\begin{equation}
\boxed{E_P=k_B\Theta_{\rm clk}\log3.}
\label{p12-83-c13-s2:eq:planck-trit}
\end{equation}

The relation does not depend on a coincidence between units: both sides
follow from the same chronological structure \(108\).

When the gravitational constant of \cref{ch:gravity} is introduced, the Planck family satisfies

\begin{align}
\ell_P^2&=\frac{G\hbar}{c^3}=\theta_{108}\ell_0^2,
\label{p12-83-c13-s2:eq:planck-area}\\
F_P&=\frac{c^4}{G},
\qquad
P_P=\frac{c^5}{G},
\label{p12-83-c13-s2:eq:planck-force-power}\\
\rho_P&=\frac{\hbar}{\theta_{108}^2c^5t_0^4}.
\label{p12-83-c13-s2:eq:planck-density}
\end{align}

Although \(G\) and \(\hbar\) transform reciprocally between the sections
before and after the return, their product in \(\ell_P^2\) preserves the
geometry. This invariance determines area as the interface between the
geometric, informational and entanglement realizations.

\subsection{Spectral formulations of Wien's displacement law}

For the spectral density per wavelength, the maximum satisfies

\[
5(1-\ee^{-x_5})=x_5,
\qquad
x_5=5+W_0(-5\ee^{-5}).
\]

With \cref{p12-83-c13-s2:eq:boltzmann-clock},

\begin{equation}
\lambda_{\max}(\Theta_{\rm clk})
=\frac{108\log3}{x_5}\,\ell_0.
\label{p12-83-c13-s2:eq:wien-lambda}
\end{equation}

For the frequency density,

\[
3(1-\ee^{-x_3})=x_3,
\qquad
x_3=3+W_0(-3\ee^{-3}),
\]

and

\begin{equation}
\nu_{\max}(\Theta_{\rm clk})
=\frac{x_3}{108t_0\log3}.
\label{p12-83-c13-s2:eq:wien-frequency}
\end{equation}

The maxima are not reciprocals, since the spectral densities
transform with the corresponding Jacobian. The quotient \(x_5/x_3\)
belongs to the structure of the parameterization and does not represent an
inconsistency.

\subsection{Stefan--Boltzmann law and thermodynamics of the photon gas}

Evaluating the radiation law at \(\Theta_{\rm clk}\) produces the intrinsic flux

\begin{equation}
j_\star(\Theta_{\rm clk})
=\frac{2\pi^5}{15\,108^4(\log3)^4}\,
\frac{h_{\rm int}}{\ell_0^2t_0^2},
\label{p12-83-c13-s2:eq:stefan-hmt}
\end{equation}

in the declared normalization. The Stefan--Boltzmann constant is not inserted as a primitive: it results from integrating the bosonic spectrum once action, propagation, and temperature have been fixed.
The spectral laws employed are those of Planck and Wien \cite{rm:Planck,rm:Wien}; HMT fixes here the section on which they are evaluated, not their constants by retrospective comparison.

In units of the Planck scale,

\begin{align}
n_\gamma\ell_P^3
&=\frac{2\zeta(3)}{\pi^2\log^3 3},
\label{p12-83-c13-s2:eq:photon-number}\\
\frac{u_\gamma\ell_P^3}{E_P}
&=\frac{\pi^2}{15\log^4 3},
\label{p12-83-c13-s2:eq:photon-energy}\\
\frac{s_\gamma\ell_P^3}{k_B}
&=\frac{4\pi^2}{45\log^3 3}.
\label{p12-83-c13-s2:eq:photon-entropy}
\end{align}

\(\zeta(3)\) represents the third-order bosonic moment, whereas
\(\log3\) determines the energy associated with the TRIT decision. The equation preserves the
distinction between the two invariants.

The quantum thermal conductance is expressed as

\begin{equation}
\frac{\kappa_Qt_0}{k_B}=\frac{\pi^2}{324\log3}.
\label{p12-83-c13-s2:eq:thermal-conductance}
\end{equation}

\subsection{Posterior metrological confluence of action, charge, conductance, and magnetic flux}

\begin{longtable}{@{}p{0.18\textwidth}p{0.27\textwidth}p{0.45\textwidth}@{}}
\toprule
Object & Equation internal & Significado estructural\\
\midrule
\endhead
\(R_K\) & \(Z_{\rm vac}/(2\alpha)\) & bidirectionally invariant quantum resistance\\
\(G_0\) & \(4\alpha/Z_{\rm vac}\) & conductancia complementaria\\
\(\Phi_0\) & \(\sqrt{hZ/(8\alpha)}\) & flow preserving orientation of action\\
\(K_J\) & \(\Phi_0^{-1}\) & reciprocal relation between frequency and flux\\
\(E_P\) & \(k_B\Theta\log3\) & energy of the cycle and TRIT informational content\\
Wien's law & \(108\log3/x_5\) & endpoint spectral in the scale chronological\\
Photonic Apéry & \(2\zeta(3)/(\pi^2\log^3 3)\) & number density of the bosonic field\\
\bottomrule
\end{longtable}

\begin{falsifier}
The block is falsified if a change of spectral variable improperly preserves
the same Wien extremum, if \(\zeta(3)\) replaces
\(\log3\), or if separate values are assigned to \(k_B\) and \(\Theta\) without
specifying a chart. The identities
\cref{p12-83-c13-s2:eq:quantum-electrical-identities,p12-83-c13-s2:eq:planck-trit} provide direct
algebraic controls.
\end{falsifier}

\begin{statusbox}
Electric quanta and Planck--TRIT confluence: \status{Internal exact theorem}. Wien, Stefan--Boltzmann and gas photonic: \status{Composition mathematical exact} with the declared bosonic continuum. The SI chart remains as \status{Contrast metrological external}.
\end{statusbox}


\section{Graduated analytic realization of the quintuple microfiber}
\label{p12-83-c13-s3:chap:realizacion-analitica-contraangulo}
\index{conjugate angular coordinate \(C^*\)}
\index{microfiber of pi@microfiber of \(\pi\)!analytic character}
\index{determinantal character}

The finite catalog and an analytic evaluation answer questions of different
types. The catalog determines which regions exist, which share a
return and what information each preserves. An analytic evaluation must
also decide how the length of a word is transformed into degree, how
a branch is prolonged after its return and with which scalar character that
prolongation is read. This chapter explicitly constructs that second
operation.

The construction begins with four data already obtained in the preceding
chapters: the five-region microfiber of \(\pi\), the length six of
each regional word, the dodecaphasic closure of \(\alpha\), and the bilayer
transport associated with the common coordinate
\(A=10^3\alpha\). No value of the conjugate angular coordinate \(C_*^{\rm HMT}\), of the reduced action
component, of the Barbero--Immirzi functional, or of a metrological constant is
used in the evaluation. The result is a convergent regional character,
an exact recurrence, and an oriented phase whose sexagesimal chart reproduces
the published conjugate angular coordinate.

The separation between the finite part and the analytic part is essential. The
integer \(169\) will be a theorem of the catalogue. The series that prolongs it will be the
theorem of an analytic reader whose rules are declared before it is evaluated.
Thus, the additional premises are not concealed within a decimal
coincidence.

\subsection{The persistent return of the \(5\) regions}

Let
\[
\mathcal F_\pi=\Psi^{-1}(010211)\subset\mathcal U_6
\]
the microfiber relative to the catalog. Recall its five elements, now
separated into initial block and terminal block:
\[
\begin{aligned}
 &(501,614,498\mid169,272,272),\\
 &(810,923,870\mid169,272,272),\\
 &(810,983,810\mid169,272,272),\\
 &(870,923,810\mid169,272,272),\\
 &(870,923,810\mid418,674,521).
\end{aligned}
\]
We define the return projection.
\[
p_{\rm ret}:\mathbb Z^6\longrightarrow\mathbb Z^3,
\qquad
p_{\rm ret}(u_1,\ldots,u_6)=(u_4,u_5,u_6).
\]

\begin{proposition}[Persistent terminal block]
\label{p12-83-c13-s3:prop:bloque-terminal-persistente}
The multiset \(p_{\rm ret}(\mathcal F_\pi)\) has a unique element of
maximum multiplicity:
\[
b_\pi=(169,272,272),
\qquad
\operatorname{mult}_{\mathcal F_\pi}(b_\pi)=4.
\]
The only remaining terminal block is
\[
b_\pi^-=(418,674,521)
\]
and has multiplicity one.
\end{proposition}

\begin{proof}
The assertion follows by exact enumeration of the five rows of the
catalog that satisfy \(\Psi(U)=010211\). Four rows have terminal
block \((169,272,272)\); the fifth, which carries the mutated return, has
\((418,674,521)\). The maximum multiplicity is four and is attained only
once.
\end{proof}

The oriented difference between the mutated return and the persistent is
\[
 \omega_\pi=b_\pi^- - b_\pi=(249,402,249),
 \qquad
 \langle(1,1,1),\omega_\pi\rangle=900.
\]
This identity separates two objects that must not be conflated. The integer
\(169\) is a coordinate of a terminal datum, that is, a datum of degree
zero. The vector \(\omega_\pi\) can be read as a one-cochain
after orienting the edge joining the two blocks. Calling it a cocycle would additionally
require fixing a cochain complex and verifying that its coboundary vanishes.
The analytic prolongation constructed below does not presuppose that
identification.

The selection of \(169\) does not require privileging the first position of the block. The symmetric group \(\mathfrak S_3\) acts by permuting its three coordinates. The stabilizer of \(b_\pi\) exchanges the two occurrences of \(272\) and fixes the unique coordinate of multiplicity one. Denote that coordinate by \(\operatorname{sing}(b)\) when the block has type \((r,s,s)\), with \(r\ne s\). Equivalently,
\[
\operatorname{sing}(b)
=\sum_{j=1}^3b_j-2\operatorname{moda}(b).
\]
In particular,
\[
\boxed{\operatorname{sing}(b_\pi)=169.}
\]

\begin{definition}[Residual fiber selector]
\label{p12-83-c13-s3:def:selector-residual-pi}
The selector \(\rho_\pi\) successively applies two operations:
\begin{enumerate}
  \item selects the terminal block of maximum multiplicity within
  \(p_{\rm ret}(\mathcal F_\pi)\);
  \item retains the coordinate fixed by the stabilizer of that block.
\end{enumerate}
By \cref{p12-83-c13-s3:prop:bloque-terminal-persistente},
\[
\boxed{\rho_\pi=169.}
\]
\end{definition}

\begin{theorem}[Equivariant uniqueness of return]
\label{p12-83-c13-s3:thm:unicidad-retorno-169}
Among the selectors that
\begin{enumerate}
  \item are invariant under permutations of the five regions;
  \item are equivariant under permutations of the three terminal
  coordinates;
  \item select the block of maximal multiplicity;
  \item retain a coordinate fixed by its stabilizer,
\end{enumerate}
the return value is uniquely determined and is \(169\).
\end{theorem}

\begin{proof}
Invariance under \(\mathfrak S_5\) prevents using the order of the rows. The
third condition selects \(b_\pi\), which is the unique mode by
\cref{p12-83-c13-s3:prop:bloque-terminal-persistente}. Its stabilizer in
\(\mathfrak S_3\) has two orbits on the coordinates: the pair of
values \(272\) and the singleton value \(169\). The fourth condition selects the
second orbit. Equivariance preserves its value when its position is shifted.
\end{proof}

This argument improves the positional reading of the return. The \(169\) is not
«the first coordinate» of a selected row: it is an invariant of the
regional multiset and of the action of its reordering symmetries.
Nor is it selected by global rarity. Among the \(468\) emissions in the
catalogue, \(169\) appears \(84\) times, uniformly distributed with
\(14\) occurrences in each of the six positions. The uniqueness of
\cref{p12-83-c13-s3:thm:unicidad-retorno-169} is therefore relational and regional: it resides
in the combination of microfiber, modal persistence, and stabilizer, not in the
isolated frequency of the integer.

\subsection{Cylindrical length and analytic degree}

Let \(\mathcal P\) be the free module generated by the regional words and
let \(F^n\mathcal P\) be its length filtration. The associated grading
records the first level at which each word appears. For a parameter
\(z\), the ordinary length character is
\[
\chi_z([w])=z^{|w|}.
\]
Multiplicativity
\[
\chi_z([uv])=\chi_z([u])\chi_z([v])
\]
expresses that concatenation adds lengths. When evaluated at
\(z=\alpha\), a word of length \(n\) receives degree \(\alpha^n\).

\begin{definition}[Degree compatibility]
\label{p12-83-c13-s3:def:compatibilidad-grado}
The regional analytic reader is compatible with the cylindrical filtration
when it sends the homogeneous component of length \(n\) to analytic degree
\(n\):
\[
\operatorname{gr}_n\mathcal P\longrightarrow
\alpha^n\mathbb R.
\]
\end{definition}

The five regions belong to \(\mathcal U_6\); therefore, the return datum
appears in degree six:
\[
\rho_\pi[U_6]\longmapsto169\alpha^6.
\]
This six is the length of the regional word. It is not the length of the
closed walks of \cref{chap:cinco-ciclos-pi}. Those walks traverse
the four anchors
\(U_\pi,U_e,U_\varphi,U_{\rm APP}\) simultaneously and have minimum length eight. A
word, an edge of the block graph, and a closed walk are objects of
distinct types; the grading used here belongs to the first.

\subsection{The Determinantal Character of Common Transport}

The dodecaphasic closure provides \(\alpha\). The completion of the three
nonadic pairs provides the scale \(10^3\), so that
\[
A=10^3\alpha.
\]
Having chosen the Archimedean angular chart, the common coordinate in radians is
\[
x=\frac{\pi A}{180}=\frac{50\pi}{9}\alpha.
\]
Let \(R\) be a real involution of zero trace in dimension two,
\[
R^2=I_2,
\qquad
\operatorname{tr}R=0,
\]
and consider the formal two-layer transport
\[
\mathcal T(x,y)=\exp(-xI_2-yR).
\]
The oriented variable \(y\) has not yet been evaluated. Nevertheless, the action
of \(\mathcal T\) on the determinant line is independent of it:
\[
\begin{aligned}
\det\mathcal T(x,y)
&=\exp\!\left(\operatorname{tr}(-xI_2-yR)\right)\\
&=e^{-2x}.
\end{aligned}
\]
We define
\[
\boxed{
D_A=e^{-2x}=e^{-100\pi\alpha/9}.}
\]
The subscript recalls that this is the determinant of the common contraction
associated with \(A\), not a function of \(C_*^{\rm HMT}\).

\begin{proposition}[Character of the determinant line]
\label{p12-83-c13-s3:prop:caracter-determinante}
The map
\[
\delta_A:(\mathbb N_0,+)\longrightarrow(\mathbb R_{>0},\cdot),
\qquad
\delta_A(k)=D_A^k,
\]
is the unique multiplicative character whose value on an elementary
prolongation is \(D_A\).
\end{proposition}

\begin{proof}
The multiplicativity gives
\[
\delta_A(k)
=\delta_A(\underbrace{1+\cdots+1}_{k\text{ times}})
=\delta_A(1)^k=D_A^k.
\]
The same identity proves existence and uniqueness.
\end{proof}

The determinant remains invariant under conjugation of
\(\mathcal T\) and under the exchange of sheets \(R\mapsto-R\). Therefore, the
tail constructed from \(D_A\) belongs to the common mode. The orientation
will appear in the sign with which the regional character corrects the phase.

\subsection{Rules of the Regional Analytic Reader}

Analytic continuation is fixed by the following rules.

\begin{definition}[Regional Determinantal Reader]
\label{p12-83-c13-s3:def:lector-regional-determinantal}
The determinantal regional reader satisfies:
\begin{enumerate}
  \item \emph{degree compatibility}: cylindrical length is the analytic
  degree;
  \item \emph{renewal decomposition}: the finite return impulse and
  strict continuation belong to distinct additive summands;
  \item \emph{continuation normalization}: after recording the residue
  \(\rho_\pi\) just once, the surviving branch begins at the unit of
  the determinant line;
  \item \emph{concatenation covariance}: each elementary prolongation
  acts through \(\bigwedge^2\mathcal T=D_A\);
  \item \emph{orientation}: in the cardinal convention fixed for APP, the
  regional return belongs to the compensating sector and is subtracted from the fifth-degree
action phase.
\end{enumerate}
\end{definition}

The third rule has mathematical content. The value \(169\) is recorded once as the amplitude of the return; it is not multiplied again at each prolongation. The strict continuation counts the sole line that continues and normalizes it by its unit. This is the distinction between a renewal recurrence and homogeneous iteration of the initial impulse.

The need to declare that normalization can be proved. For every
\(\lambda\in\mathbb R\), the family
\[
F_\lambda(z)
=169z^6+\lambda\frac{D_Az^7}{1-D_Az}
\]
preserves the same finite return and the same determinantal ratio between
consecutive coefficients. The catalogue and the ratio alone do not distinguish
\(\lambda\). The unit of the determinant line fixes
\(\lambda=1\) before \(z=\alpha\) is evaluated. Thus, the normalization
is not obtained by fitting the final value; it forms part of the definition of the
character.

\subsection{Return recurrence and closed sum}

We define the coefficients \(\kappa_n\) by
\[
\kappa_6=169,
\qquad
\kappa_7=D_A,
\qquad
\kappa_{n+1}=D_A\kappa_n\quad(n\ge7).
\]
The first equality is the regional impulse; the next two describe the normalized continuation.

\begin{theorem}[Graduated analytic realization]
\label{p12-83-c13-s3:thm:realizacion-analitica-graduada}
The analytic germ compatible with the rules of
\cref{p12-83-c13-s3:def:lector-regional-determinantal} is unique and equals
\[
\boxed{
\mathscr C_\pi(z)
=169z^6+\frac{D_Az^7}{1-D_Az}.}
\]
Their coefficients satisfy
\[
\kappa_n=D_A^{n-6}\qquad(n\ge7).
\]
\end{theorem}

\begin{proof}
Write
\[
F(z)=169z^6+\sum_{k\ge1}c_kz^{6+k}.
\]
The normalization of continuation and the determinantal covariance give
\[
c_1=D_A,
\qquad
c_{k+1}=D_Ac_k.
\]
By induction, \(c_k=D_A^k\). The geometric sum yields
\[
\sum_{k\ge1}D_A^kz^{6+k}
=\frac{D_Az^7}{1-D_Az}.
\]
No free coefficients remain.
\end{proof}

Since \(0<\alpha<1\) and
\(0<D_A<1\),
\[
0<D_A\alpha<1.
\]
Consequently, \(\mathscr C_\pi(\alpha)\) converges absolutely and its
radius of convergence is \(D_A^{-1}>1\). Numerically,
\[
D_A=0.7751291192471564301888840964802\ldots,
\]
\[
D_A\alpha=0.00565639046986492673885079095469\ldots.
\]
The speed of this contraction explains why the complete sum and its
degree-seven truncation have the same first fifteen digits in the
final chart.

\subsection{The 5th-degree action character}

The central block
\[
\mathcal B_c=\begin{pmatrix}7&2\\2&7\end{pmatrix}
\]
has eigenvalues \(9\) and \(5\). The dodecaphasic cycle contributes rank
\(12\); its step-three orbits have length four; and the complete census
contributes the exact quotient
\[
\frac{104976}{1944}=54.
\]
From these invariants the fifth-degree action character is fixed.
\[
\boxed{
H_5(\alpha,\varphi)
=\frac12\sqrt{\frac{1000\alpha}{\varphi}}-\alpha
+\frac9{16}\alpha^2-\frac59\alpha^3
+\frac7{48}\alpha^4-\frac1{54}\alpha^5.}
\]
The four rational coefficients can be written as
\[
\frac9{16}=\frac{\lambda_+}{4^2},
\qquad
\frac59=\frac{\lambda_-}{\lambda_+},
\qquad
\frac7{48}=\frac{\tfrac12\operatorname{tr}\mathcal B_c}{4\cdot12},
\qquad
\frac1{54}=\frac{1944}{104976},
\]
with \((\lambda_+,\lambda_-)=(9,5)\). These identities certify the
provenance of the numbers used. The successive assignment of those
invariants to degrees two, three, four, and five forms part of the analytic
character; it does not follow from the mere existence of the finite catalog.

This precision removes a common ambiguity. A combination of
invariants of the construction is an internal and reproducible definition; its
naturality requires the reader rules specifying order, signs, and
normalization. Here those rules are visible and are subject to the same
negative controls as the rest of the construction.

\subsection{Oriented phase and conjugate angular coordinate}

The fifth-degree phase prior to return is
\[
y_5=\frac\pi{90}\bigl(H_5+\alpha\bigr).
\]
The complete regional character produces the correction
\(\mathscr C_\pi(\alpha)\). With the orientation fixed in
\cref{p12-83-c13-s3:def:lector-regional-determinantal}, we define
\[
\boxed{
y_*=\frac\pi{90}\bigl(H_5+\alpha\bigr)
-169\alpha^6
-\frac{D_A\alpha^7}{1-D_A\alpha}.}
\]
This is the primary formula. The sexagesimal unit does not come into play until applying the chart.
\[
C_*^{\rm HMT}=\frac{180}{\pi}y_*.
\]

\begin{theorem}[Conjugate angular coordinate of the regional reader determinantal]
\label{p12-83-c13-s3:thm:contraangulo-regional}
With the dodecaphasic closure of \(\alpha\),
\(\varphi=(1+\sqrt5)/2\) and the rules of the determinantal regional reader fixed, the
phase \(y_*\) and its conjugate angular coordinate are determined without introducing the sought
value. Their evaluation is
\[
\boxed{
y_*=0.03706622646583826398493779409230638\ldots,}
\]
\[
\boxed{
C_*^{\rm HMT}
=2.12373833896864572426195138039336\ldots{}^\circ.}
\]
Therefore, to fifteen decimal places,
\[
\boxed{C_*^{\rm HMT}=2.123738338968646^\circ.}
\]
\end{theorem}

\begin{proof}
\cref{p12-83-c13-s3:thm:unicidad-retorno-169} determines \(169\). The closure of
\(\alpha\) determines \(A\), \(x\), and
\(D_A=e^{-100\pi\alpha/9}\). \cref{p12-83-c13-s3:thm:realizacion-analitica-graduada}
determines the complete series. The formula for \(H_5\) contains only
\(\alpha\), \(\varphi\), and declared structural invariants. Substituting
these data into the expression for \(y_*\) gives the printed value. The decimal value
of \(C_*^{\rm HMT}\) occurs in none of the right-hand sides.
\end{proof}

The reduced action component after return is now as follows
output:
\[
\boxed{
\bar h_{\rm HMT}^{\rm ret}
=\frac{C_*^{\rm HMT}}2-\alpha
=H_5-\frac{90}{\pi}\mathscr C_\pi(\alpha).}
\]
Numerically,
\[
\bar h_{\rm HMT}^{\rm ret}
=1.05457181691503906113369058472430\ldots.
\]
The inverse identity has not been used to introduce \(C_*^{\rm HMT}\); it is
obtained after generating \(y_*\).

\subsection{Degree-\(7\) truncation and exact remainder bound}

The expression initially found,
\[
y_7=\frac\pi{90}(H_5+\alpha)-169\alpha^6-D_A\alpha^7,
\]
is the seventh-degree truncation of the full character. Produce
\[
C_7=2.12373833896864600265403827103919\ldots{}^\circ.
\]
The omitted remainder is not an indeterminate term; it is the exact geometric tail.
\[
\mathscr C_\pi(\alpha)
-169\alpha^6-D_A\alpha^7
=\frac{D_A^2\alpha^8}{1-D_A\alpha}.
\]
Therefore,
\[
\left|C_7-C_*^{\rm HMT}\right|
=\frac{180}{\pi}\frac{D_A^2\alpha^8}{1-D_A\alpha}
=2.7839208689064583\ldots\times10^{-16}\;{}^\circ.
\]
The recurrence thus transforms a seventh-degree approximation into a
complete analytic realization and provides a closed-form bound for each
subsequent truncation.

\subsection{Ablations of the selector and continuation}

The construction changes when any of its active components is removed.
The following values are obtained without recalibrating the other terms:
\[
\begin{array}{@{}lc@{}}
\toprule
\text{variant}&\text{conjugate angular coordinate in degrees}\\
\midrule
\rho_\pi=168&2.1237383389772976863904487\ldots\\
\rho_\pi=169&2.1237383389686457242619514\ldots\\
\rho_\pi=170&2.1237383389599937621334541\ldots\\
\text{without determinantal tail}&2.1237383389686949415301675\ldots\\
D_A\text{ replaced by }1&2.1237383389686313409965826\ldots\\
\bottomrule
\end{array}
\]
These ablations do not in themselves constitute a probabilistic estimate.
They demonstrate that the equivariant selector and the determinantal character are
effective components of the output and cannot be removed while maintaining the
same realization.

\subsection{Spectral reconstruction and hexad-octad transition}

The common phase and the oriented phase determine the two channels.
\[
q_-=e^{-x+y_*},
\qquad
q_+=e^{-x-y_*}.
\]
The invariants are recovered exactly.
\[
q_-q_+=D_A,
\qquad
\alpha=-\frac9{100\pi}\log D_A,
\qquad
C_*^{\rm HMT}=\frac{90}{\pi}\log\frac{q_-}{q_+}.
\]
The first equality reads the common contraction; the third, the oriented
separation of the sheets. The exchange \(R\mapsto-R\) permutes
\(q_+\) and \(q_-\), preserves \(D_A\) and reverses the sign of \(y_*\).

For
\[
S_{90}(q)=\frac{q^{90}}{1-q^{270}},
\qquad
S_{120}(q)=\frac{q^{120}}{1-q^{360}},
\]
the hexad--octad transition functional is written entirely in the
phase chart:
\[
\Gamma_{6\to8}
=\frac{180}{\pi}xy_*
+12\bigl(S_{90}(q_-)-S_{90}(q_+)\bigr)
-\bigl(S_{120}(q_-)-S_{120}(q_+)\bigr).
\]
The evaluation produced by the coordinate \(C_*^{\rm HMT}\) of
\cref{p12-83-c13-s3:thm:contraangulo-regional} is
\[
\boxed{
\Gamma_{6\to8}
=0.274007672694459274747255260774\ldots.}
\]
Thus, the internal value recognized in Dimensional Mechanics as the Barbero--Immirzi functional ceases to receive \(C_*^{\rm HMT}\) as an independent literal: both descend from the same pair of phases generated in this chapter. The geometric and physical interpretation of \(\Gamma_{6\to8}\) belongs to the work on Dimensional Mechanics; its mathematical evaluation already belongs to the previously constructed interface.

\subsection{Action prior and posterior to return}

The fifth-degree character and the post-action component after return
are two distinct values:
\[
H_5
=1.05457181764615446962582182943306\ldots,
\]
\[
\bar h_{\rm HMT}^{\rm ret}
=1.05457181691503906113369058472430\ldots.
\]
The first is the local action before incorporating the regional character; the
second is the oriented action remaining after the return. The correction
that generates the conjugate angular coordinate prevents identifying them.

The metrological comparison must respect this distinction. The first value gives
the mantissa
\[
2\pi H_5
=6.62607014999998792600111034989947\ldots,
\]
while the subsequent action gives
\[
2\pi\bar h_{\rm HMT}^{\rm ret}
=6.62607014540625433351074984759288\ldots.
\]
In the International System, the value
\[
h=6.62607015\times10^{-34}\;\mathrm{J\,s}
\]
is exact by definition \cite{BIPMSI}. By therefore,
\[
2\pi H_5-6.62607015
=-1.2073998889650101\ldots\times10^{-14},
\]
and it is not an exact equality. The proximity of the value before the return is
a remarkable numerical check; it does not authorize replacing the defining constant
of the SI with the subsequent component or confusing the two HMT actions\@.

This result precisely determines the mathematical situation. The regional
reader generates \(C_*^{\rm HMT}\) and, from it, the functional
\(\Gamma_{6\to8}\). The same calculation separates the Planck approximation
produced by \(H_5\) from the reduced action intervening in
the conjugate angular coordinate. The separation does not weaken the construction: it eliminates an
identification incompatible with the exact figures and converts every
comparison into a falsifiable assertion of a well-defined type.

\subsection{Relative transporter and spectral disjunction of the blocks}
\label{p12-83-c13-s3:sec:transportador-relativo-split}

In the ordered basis of channels, let
\[
 R=\begin{pmatrix}0&1\\1&0\end{pmatrix},
 \qquad P_\pm=\frac{I\pm R}{2},
 \qquad
 B_0=7I+2R,
 \qquad B_1=\Adeg I+\Cdeg R.
\]
The split decomposition
\[
 xI+yR=\rho(x,y)e^{\eta(x,y)R},
 \quad \rho(x,y)=\sqrt{x^2-y^2},
 \quad \eta(x,y)=\operatorname{artanh}(y/x)
\]
constructs the unique transporter
\[
 T_{\rm rel}
 =\frac{\sqrt{(\Adeg)^2-(\Cdeg)^2}}{\sqrt{45}}
   \exp\!\left[
   \left(\operatorname{artanh}\frac{\Cdeg}{\Adeg}
   -\operatorname{artanh}\frac27\right)R\right],
 \qquad T_{\rm rel}B_0=B_1.
\]
In the spectral basis it acts by
\(9\mapsto\Adeg+\Cdeg\) and
\(5\mapsto\Adeg-\Cdeg\).

\begin{theorem}[Spectral Disjunction of the Local and Angular Realizations]
\label{p12-83-c13-s3:thm:no-go-entrelazador-bloques-split}
The spectra
\[
 \operatorname{Spec}(B_0)=\{9,5\},\qquad
 \operatorname{Spec}(B_1)=\{\Adeg+\Cdeg,\Adeg-\Cdeg\}
\]
are disjoint. Consequently,
\[
 \operatorname{Hom}_{B_0,B_1}
 :=\{\Phi:\Phi B_0=B_1\Phi\}=\{0\}.
\]
\end{theorem}

\begin{proof}
The rational bounds constructed for the angular coordinates give
\(7.29<\Adeg<7.30\) and \(2.12<\Cdeg<2.13\); therefore,
\(9.41<\Adeg+\Cdeg<9.43\) and
\(5.16<\Adeg-\Cdeg<5.18\). In the basis that diagonalizes \(R\), each entry
of an intertwiner is multiplied by the difference between an eigenvalue
of \(B_0\) and one of \(B_1\); disjointness forces all of them to vanish.
\end{proof}

The relative transporter compares two blocks within the split algebra;
the dodecaphasic--Witt and incidence intertwiners act between
equivalent representations in their own spaces of residence. This separation
positively types both operations and preserves their domains.

\subsection{Generation chain \(\mathcal F_\pi\to b_\pi\to\rho_\pi\to D_A\to C_*^{\rm HMT}\to\Gamma_{6\to8}\): regional selection, determinantal prolongation and oriented coordinate}

The entire procedure can be summarized as a chain of applications:
\[
\begin{aligned}
\mathcal F_\pi
&\longmapsto b_\pi
\longmapsto\rho_\pi=169,\\
\alpha
&\longmapsto A
\longmapsto x
\longmapsto D_A,\\
(\rho_\pi,D_A)
&\longmapsto\mathscr C_\pi(\alpha),\\
(H_5,\mathscr C_\pi(\alpha))
&\longmapsto y_*
\longmapsto C_*^{\rm HMT},\\
(x,y_*)
&\longmapsto(q_+,q_-)
\longmapsto\Gamma_{6\to8}.
\end{aligned}
\]
The first line belongs to the finite catalog; the second, to the common closure of
\(\alpha\); the third, to the determinantal prolongation; the fourth produces the
oriented coordinate; the fifth transports it toward the dimensional interface.

In this organization, \(\alpha\) measures the common contraction of the transport, and \(C_*^{\rm HMT}\) measures the oriented separation of its two sheets. The regional return adds no external constant: it selects a finite amplitude and prolongs it through the determinant already fixed by \(\alpha\). The chart in degrees is subsequent. The primary object is the pair of phases \((x,y_*)\).

An exhaustive verification reconstructs the five regions directly
from the catalogue, traverses the actions of
\(\mathfrak S_5\) and \(\mathfrak S_3\), checks the recurrence, and sums the
series. The input of \(\alpha\) is the exact output of its
dodecaphasic closure, not a transcribed metrological digit. The verification keeps
the catalogue theorems, reader rules, and external comparison
separate; therefore, no experimental comparison intervenes in the generating
function.
